\documentclass[11pt,a4paper]{article}
\usepackage[T1]{fontenc}
\usepackage{mathptmx}
\usepackage[margin=25mm]{geometry}
\usepackage{booktabs}
\usepackage{array}
\usepackage{amsmath}
\usepackage{microtype}
\usepackage{titlesec}
\usepackage[hidelinks]{hyperref}
\usepackage{fancyhdr}
\titleformat{\section}{\large\bfseries}{\thesection}{0.6em}{}
\titleformat{\subsection}{\bfseries}{\thesubsection}{0.6em}{}
\setlength{\parskip}{0.5em}
\setlength{\parindent}{0pt}
\pagestyle{fancy}\fancyhf{}
\renewcommand{\headrulewidth}{0.4pt}
\newcommand{\inr}[1]{INR~#1}

\fancyhead[L]{RescueSwarm --- Track A: Air Vehicle}
\fancyhead[R]{\thepage}
\begin{document}

\begin{center}
{\LARGE\bfseries Track A --- Air Vehicle}\\[3mm]
{\large RescueSwarm: work package and funding request}\\[2mm]
NIDAR 2026--27 $\cdot$ Track 1 $\cdot$ Mission 1\\[1mm]
\textbf{Ask: \inr{325,397}} $\cdot$ 2 students $\cdot$ Frame, propulsion, power, payload mechanism, assembly
\end{center}
\vspace{2mm}\hrule\vspace{4mm}


This track builds the thing that flies. Everything else in the programme is
carried by it, which means this track's schedule is the programme's schedule:
no other track can begin flight validation until an aircraft exists, and the
single measurement that decides whether one does is owned here.

\section{Scope and ownership}

This track owns \textbf{frame, propulsion, power, payload mechanism, assembly}. It is one of four delivery tracks defined in
\texttt{docs/development-plan.md}~\S1.3; the integrated system argument lives in
the master proposal, \texttt{docs/proposal/rescueswarm-proposal.pdf}, which this
document does not restate.

\section{The design point this track is funded to build}

The aircraft is a quadrotor on 18\,in propellers with a maximum take-off mass of
\textbf{6.36\,kg}, giving a three-aircraft fleet of \textbf{19.08\,kg} against
the 25\,kg regulatory cap --- 24\,\% margin. Structure is 1\,495\,g (23.5\,\% of
MTOW) and the battery pack 1\,449\,g (22.8\,\%); together they are 46\,\% of the
aircraft, which is why airframe mass fraction and pack chemistry are the two
sensitivities this track carries in the risk register.

All of these figures regenerate from \texttt{tools/sizing-model/}\allowbreak
\texttt{rescueswarm\_sizing\_model.py}. None is transcribed by hand.

\section{Propulsion, and the measurement that is load-bearing}

At a thrust-to-weight ratio of 2.0 the aircraft requires \textbf{125\,N total},
which is \textbf{3.18\,kgf per motor} on an 18\,in propeller, with hover sitting
at 1.59\,kgf --- 50\,\% of maximum, inside the healthy 45--55\,\% band. The mass
budget allows 160\,g per motor.

\textbf{The funded motors publish no thrust curve.} That is a deliberate cost
decision, not an oversight: professional-grade motors with published data cost
roughly 2.4 times as much. The consequence is that the thrust-to-weight figure
above is a \emph{requirement} rather than a measurement until the thrust stand
runs in P2, and the thrust stand is therefore not a convenience --- it is the
instrument that says whether this aircraft hovers. It is funded accordingly.

Two figures in \texttt{docs/sizing/sizing-calculations.md} contradict the model
here, quoting 2.53 and 2.94\,kgf per motor against the model's 3.18, and one
line still specifies a 20\,in propeller where the bill of materials buys 18\,in.
Reconciling those is a P1 action for this track.

\section{Energy: what the pack must actually do}

The pack is sized by a reserve policy, not by the design mission. The mission
consumes 105.5\,Wh; the policy requires the pack to also carry one complete
re-sweep and four minutes of loiter, all inside an 80\,\% depth of discharge:
\[
105.5 + 22.1 + 60.8 = 188.4~\text{Wh usable},
\]
against 233.3\,Wh available --- a \textbf{24\,\% margin}.

Stated independently of chemistry, so that any candidate pack can be tested
against it, the requirement is \textbf{6S}, \textbf{$\geq$236\,Wh} nameplate,
\textbf{$\geq$115\,A continuous}, \textbf{$\leq$40\,m$\Omega$} and
\textbf{$\leq$1\,450\,g}. The current figure is the one most often missed: it is
not a burst rating. At T/W 2.0 the aircraft draws 2\,481\,W, which at 21.6\,V is
115\,A, and that must be available continuously.

The adopted implementation is 18 cells in 6S3P at 4500\,mAh and 45\,A
continuous. A 40\,A cell leaves 4.5\,\% burst margin against 18\,\%, and cells
whose 45\,A rating depends on an 80\,\textdegree C cut-off are 35\,A cells in
Indian ambient --- below the 38.3\,A per-cell peak. The specification was
tightened accordingly on 2026-08-20.

\section{Two open items this track owns}

\textbf{The mass statement does not close.} Itemised masses sum to 6\,061\,g
against a 6\,360\,g MTOW, leaving \textbf{299\,g} --- 4.7\,\% --- unattributed.
It is small enough not to threaten fleet margin and large enough that it should
be attributed rather than absorbed. Closing it is a P1 action.

\textbf{Cell sourcing is a live commercial question.} The budget carries
\inr{700} per cell against a domestic supplier quote. A drop-in equivalent lists
publicly at \inr{405}. The action is not to switch --- cells are one of only
four remaining Indian lines and indigenisation is scored --- but to hold the
supplier quote against a published benchmark before committing.

\section{Funding request}

Every figure below is generated from
\texttt{docs/proposal/figures/track\_budget.py}, which partitions the master
competition budget. The four track asks plus the shared pool reconcile to the
master figure exactly; that reconciliation is asserted in code and fails the
build if it ever stops holding. \textbf{K} marks a line held at professional
grade on purpose.


\begin{center}
\begin{tabular}{@{}p{88mm}cr@{}}
\toprule
\textbf{Item} & & \textbf{\inr{}} \\
\midrule
Motors (4 x 4,500 x 3 ac) &  & 54,000 \\
Structure (1 x 13,000 x 3 ac) & \textbf{K} & 39,000 \\
Li-ion cells (18 x 700 x 3 ac) & \textbf{K} & 37,800 \\
Pack, BMS, PDB, BEC (1 x 8,500 x 3 ac) &  & 25,500 \\
ESCs (4 x 1,800 x 3 ac) &  & 21,600 \\
Battery chargers &  & 14,000 \\
Payload system (1 x 4,500 x 3 ac) &  & 13,500 \\
Propellers (4 x 1,000 x 3 ac) &  & 12,000 \\
Calibrated scale + remaining &  & 12,000 \\
Thrust stand &  & 8,000 \\
Wiring, connectors (1 x 2,400 x 3 ac) &  & 7,200 \\
Prop adapters (4 x 350 x 3 ac) &  & 4,200 \\
\midrule
Subtotal & & 248,800 \\
Customs duty and freight & & 13,949 \\
GST & & 20,205 \\
Contingency at 15\% & & 42,443 \\
\midrule
\textbf{TRACK A ASK} & & \textbf{325,397} \\
\bottomrule
\end{tabular}
\end{center}


\textbf{Deferred or institutionally held, and therefore not requested:}
3D printer, Cell tester, Charger PSU, Recovery parachutes, 3, Relief kits (14), Rotary tool + CF extraction, Spare airframe structure set, Spare battery packs, Spare motors, Spare propellers, plate and tube stock. These remain capabilities the track requires; they are simply not being
bought. Where a deferral carries a technical consequence it is stated in the
risk table below rather than left implicit.

\section{Deliverables by phase}

\begin{center}
\begin{tabular}{@{}llp{88mm}@{}}
\toprule
\textbf{Phase} & \textbf{Dates} & \textbf{This track delivers} \\
\midrule
\textbf{P1} & 24 Aug -- 13 Sep & Long-lead orders placed: cells, motors, structure stock. Mass residual attributed. Sizing-document contradictions reconciled. \\
\textbf{P2} & 14 Sep -- 4 Oct & \textbf{Thrust-stand measurement} -- the gate for the whole programme. Pack assembled and bench-discharged. \\
\textbf{P3} & 5 Oct -- 18 Oct & First flight. Measured mass and hover current reported to Track C. \\
\textbf{P5} & 9 Nov -- 22 Nov & Single-aircraft full mission including payload release. \\
\bottomrule
\end{tabular}
\end{center}

\section{Risks owned by this track}

\begin{center}
\begin{tabular}{@{}p{50mm}lp{70mm}@{}}
\toprule
\textbf{Risk} & \textbf{Sev.} & \textbf{Mitigation} \\
\midrule
Motors deliver less than 3.18\,kgf & High & P2 thrust stand. If they miss, the fallback is a published-curve motor at roughly 2.4$\times$ the cost, which is a budget event this track must flag the moment it is known. \\
Mass residual grows during build & Medium & 299\,g is unattributed today. Fleet carries 4\,919\,g of growth allowance, so this is affordable, but it must be tracked from P1 not discovered at weigh-in. \\
Cell lead time & Medium & No spare packs are funded. Domestic sourcing is 2--3 weeks; order at P1, not P2. \\
\bottomrule
\end{tabular}
\end{center}

\section{Interfaces to other tracks}

\begin{center}
\begin{tabular}{@{}lp{110mm}@{}}
\toprule
\textbf{With} & \textbf{Dependency} \\
\midrule
Track B & Provides the power distribution the avionics run on, and the vibration-isolated mounting the flight controller requires. \\
Track C & Provides the mass, thrust and endurance figures the SITL model flies against. If the built aircraft differs from the model, the simulations stop being evidence. \\
Track D & Provides the camera mounting plane and its boresight stability. Geolocation accuracy depends on that mount not moving. \\
\bottomrule
\end{tabular}
\end{center}

Interfaces are where student programmes fail, because each side assumes the
other owns the gap. Every row above is a two-way commitment and should be
recorded as an interface control document at P1.


\vfill
\hrule\vspace{2mm}
{\small Generated by \texttt{tools/proposal/build\_track\_proposals.py} from
\texttt{docs/proposal/figures/track\_budget.py}. Financial figures are not
maintained in this document and must not be edited here: change the master
budget and regenerate. Technical claims cite the model or document that
produced them.}
\end{document}
